% =====================================================================
%  Paper de Workshop - Avances de Tesina (v2, QSVM como modelo central)
%  Compilable en Overleaf (formato IEEE conference)
% =====================================================================
\documentclass[10pt,conference]{IEEEtran}
\IEEEoverridecommandlockouts

\usepackage[utf8]{inputenc}
\usepackage[spanish,es-noshorthands]{babel}
\usepackage{cite}
\usepackage{amsmath,amssymb,amsfonts}
\usepackage{graphicx}
\usepackage{xcolor}
\usepackage{booktabs}
\usepackage{multirow}
\usepackage{hyperref}
\usepackage{tikz}
\usetikzlibrary{positioning,arrows.meta,shapes.geometric,calc,backgrounds,fit}

\hypersetup{colorlinks=true,linkcolor=blue,citecolor=blue,urlcolor=blue}

% Traducciones de IEEEtran al español
\renewcommand{\IEEEkeywordsname}{Palabras clave}
\renewcommand{\IEEEproofname}{Demostración}

% Comando auxiliar para valores pendientes de completar
\newcommand{\pendiente}[1]{\textcolor{blue}{\textbf{[#1]}}}

\begin{document}

% ============================ PORTADA =================================
\title{Detección de Anomalías en Tráfico de Red mediante Máquinas de
Vectores de Soporte Cuánticas:\\ Validación Empírica en Hardware Cuántico
de Resonancia Magnética Nuclear}

\author{
\IEEEauthorblockN{Ticiana Angelucci}
\IEEEauthorblockA{Licenciatura en Sistemas\\
Universidad Champagnat\\
Mendoza, Argentina}
\and
\IEEEauthorblockN{Rodrigo Elgueta}
\IEEEauthorblockA{Director de Tesina\\
Universidad Champagnat\\
Mendoza, Argentina}
\and
\IEEEauthorblockN{Andrés A. Reynoso}
\IEEEauthorblockA{Asesor Científico\\
Instituto Balseiro\\
Bariloche, Argentina}
}

\maketitle

% ============================ RESUMEN =================================
\begin{abstract}
Los Sistemas de Detección de Intrusos (IDS) basados en Machine Learning
clásico enfrentan limitaciones crecientes ante el volumen y la
complejidad del tráfico de red moderno. El Quantum Machine Learning
(QML) ofrece un paradigma alternativo que proyecta los datos en espacios
de características de alta dimensionalidad (espacios de Hilbert). Este
trabajo presenta una Máquina de Vectores de Soporte Cuántica (QSVM)
basada en un kernel cuántico de fidelidad para la clasificación binaria
de tráfico de red benigno y malicioso, empleando el dataset
CIC-IDS-2017. Se documenta la decisión de diseño de abandonar un
Clasificador Cuántico Variacional (VQC) --afectado por el problema de
mesetas estériles (\textit{barren plateaus})-- en favor del enfoque de
kernel, más robusto para datos tabulares de ciberseguridad. El modelo
QSVM alcanzó una exactitud del $83{,}3\%$ en simulación clásica. La
validación se trasladó a hardware cuántico real de Resonancia Magnética
Nuclear (SpinQ Triangulum, 3 qubits), donde una prueba sobre tráfico en
vivo arrojó una exactitud del $75\%$, una precisión del $100\%$ y un
F1-Score del $80\%$, con una desviación media del kernel del
$12{,}94\%$ respecto del valor ideal. Los resultados evidencian tanto
la viabilidad funcional del enfoque como el impacto del ruido inherente
a la era NISQ.
\end{abstract}

\begin{IEEEkeywords}
Quantum Machine Learning, QSVM, kernel cuántico, detección de intrusos,
era NISQ, resonancia magnética nuclear, CIC-IDS-2017, ciberseguridad.
\end{IEEEkeywords}

% ============================ I. INTRODUCCIÓN =========================
\section{Introducción}
La evolución y sofisticación de los ciberataques, sumada al crecimiento
exponencial del tráfico en las infraestructuras de red, ha llevado a los
Sistemas de Detección de Intrusos (IDS) clásicos a un punto de
saturación. Los enfoques basados en firmas y, más recientemente, en
Machine Learning (ML) clásico, alcanzan un techo de rendimiento y
escalabilidad frente a la alta dimensionalidad de los datos de red
modernos, lo que se traduce en tasas elevadas de falsos positivos y
negativos.

El Quantum Machine Learning (QML) emerge como un cambio de paradigma
prometedor \cite{biamonte2017}. Su ventaja teórica radica en la
capacidad de mapear datos complejos en espacios de Hilbert de
dimensionalidad exponencial, habilitando fronteras de separación
inaccesibles para los métodos clásicos. No obstante, la industria
transita actualmente la era NISQ (\textit{Noisy Intermediate-Scale
Quantum}) \cite{preskill2018}, caracterizada por procesadores de escala
intermedia con alta susceptibilidad al ruido, la decoherencia y los
errores de compuerta.

Esta investigación se centra en una Máquina de Vectores de Soporte
Cuántica (QSVM) que emplea un circuito cuántico para estimar un kernel
de similitud entre muestras de tráfico de red, delegando la
clasificación final en una SVM clásica. A diferencia de los circuitos
variacionales puros, este enfoque evita los bucles de optimización
cuántica iterativa y, con ello, el colapso por mesetas estériles.

Las principales contribuciones de este reporte de avance son:
\begin{itemize}
  \item El diseño e implementación de una QSVM basada en un kernel
  cuántico de fidelidad para la clasificación binaria de tráfico de
  red, con un \textit{pipeline} de preprocesamiento adaptado a un
  procesador de 3 qubits.
  \item El análisis comparativo que motivó la transición desde un
  Clasificador Cuántico Variacional (VQC) hacia el enfoque de kernel,
  documentando el fenómeno de mesetas estériles sobre datos tabulares.
  \item Resultados empíricos de ejecución en un procesador cuántico de
  Resonancia Magnética Nuclear (SpinQ Triangulum), contrastados con la
  simulación clásica ideal.
\end{itemize}

El resto del documento se organiza de la siguiente manera: la Sección
II presenta el marco teórico; la Sección III detalla la metodología; la
Sección IV describe la implementación; la Sección V expone los
resultados experimentales; la Sección VI discute los hallazgos; y la
Sección VII concluye y delinea el trabajo futuro.

% ============================ II. MARCO TEÓRICO =======================
\section{Marco Teórico y Trabajos Relacionados}

\subsection{IDS y Machine Learning clásico}
Históricamente, la detección de intrusos ha dependido de firmas
conocidas o de anomalías estadísticas simples. La literatura reciente
documenta una transición hacia algoritmos de ML clásico
(\textit{Random Forest}, SVM, redes neuronales profundas), que mejoran
la clasificación pero enfrentan barreras de escalabilidad. En este
trabajo se adopta un modelo \textit{Random Forest} como línea base
(\textit{baseline}) de referencia.

\subsection{Quantum Machine Learning y kernels cuánticos}
El QML explota la superposición y el entrelazamiento para procesar
información \cite{biamonte2017}. Dentro de este campo, los métodos
basados en kernels cuánticos resultan particularmente atractivos para la
era NISQ: un circuito cuántico actúa como un mapa de características
(\textit{feature map}) que proyecta los datos clásicos en un espacio de
Hilbert de alta dimensionalidad, y la similitud entre muestras se
estimación mediante una medida de fidelidad cuántica. La clasificación
final recae en un modelo clásico (SVM), lo que confiere estabilidad al
conjunto.

\subsection{Clasificadores variacionales y mesetas estériles}
Los Clasificadores Cuánticos Variacionales (VQC) son el análogo
cuántico de una red neuronal: combinan un \textit{feature map} con un
circuito parametrizado (\textit{ansatz}) cuyos pesos se optimizan
iterativamente con un optimizador clásico (p.~ej., COBYLA). Sin
embargo, los VQC sufren de \textit{barren plateaus}: a medida que
crecen los qubits y la profundidad del circuito, los gradientes se
aplanan exponencialmente hacia cero y el optimizador queda «ciego».
Este fenómeno, sumado a la incompatibilidad de los datos tabulares de
red con la geometría natural de los circuitos profundos, motivó la
adopción del enfoque QSVM en este trabajo (Sección III-D).

\subsection{Dataset CIC-IDS-2017}
Se emplea el conjunto de datos público CIC-IDS-2017
\cite{sharafaldin2018}, generado por el \textit{Canadian Institute for
Cybersecurity}. Contiene tráfico benigno y múltiples tipos de ataques
(DDoS, \textit{botnet}, fuerza bruta), con más de 70 características
por conexión y etiquetado supervisado. Su uso garantiza validez
científica y comparabilidad con el estado del arte.

\subsection{Hardware cuántico de Resonancia Magnética Nuclear}
La validación física se realiza sobre un procesador SpinQ Triangulum,
una computadora cuántica de escritorio basada en Resonancia Magnética
Nuclear (RMN) de 3 qubits. A diferencia de las arquitecturas
superconductoras, la RMN opera mediante pulsos de radiofrecuencia sobre
un conjunto macroscópico de moléculas y devuelve distribuciones de
conteos por cada ejecución. Sus principales limitaciones son los tiempos
de relajación ($T_1$) y de decoherencia ($T_2$), que degradan los
estados cuánticos durante la ejecución del circuito.

% ============================ III. METODOLOGÍA ========================
\section{Metodología}

La investigación adopta un enfoque cuantitativo y experimental,
estructurado en el \textit{pipeline} ilustrado en la
Fig.~\ref{fig:arquitectura}.

\begin{figure}[!t]
\centering
\begin{tikzpicture}[
  node distance=0.75cm and 0.5cm,
  box/.style={draw, rounded corners=2pt, align=center, font=\scriptsize,
              minimum height=0.85cm, minimum width=2.1cm, fill=blue!8},
  qbox/.style={draw, rounded corners=2pt, align=center, font=\scriptsize,
              minimum height=0.85cm, minimum width=2.1cm, fill=orange!15},
  arr/.style={-{Stealth[length=2.5mm]}, thick}
]
\node[box] (data) {Dataset\\CIC-IDS-2017};
\node[box, right=of data] (pre) {Preprocesado\\SelectKBest\\$+$ Escalado $[-1,1]$};
\node[qbox, right=of pre] (fm) {Feature Map\\ZZFeatureMap};
\node[qbox, right=of fm] (kernel) {Kernel de\\Fidelidad};
\node[box, right=of kernel] (svm) {SVM\\Clásica};
\node[box, below=0.9cm of svm] (out) {Clasificación\\Benigno / Ataque};

\node[box, below=0.9cm of fm, fill=green!10] (sim) {Simulador\\clásico (ideal)};
\node[box, right=0.6cm of sim, fill=red!10] (hw) {Hardware real\\SpinQ (ruido)};

\draw[arr] (data) -- (pre);
\draw[arr] (pre) -- (fm);
\draw[arr] (fm) -- (kernel);
\draw[arr] (kernel) -- (svm);
\draw[arr] (svm) -- (out);
\draw[arr] (fm) -- (sim);
\draw[arr] (fm) -- (hw);
\draw[arr] (sim) -- (kernel);
\draw[arr] (hw) -- (kernel);

\begin{scope}[on background layer]
\node[draw, dashed, rounded corners=3pt, inner sep=4pt,
      fit=(fm)(kernel)(sim)(hw), fill=gray!5] {};
\end{scope}
\node[font=\scriptsize\itshape, anchor=north west] at
  ($(fm.north west)+(0,0.65)$) {Núcleo cuántico (Qiskit)};
\end{tikzpicture}
\caption{Arquitectura general del \textit{pipeline} QSVM. El núcleo
cuántico puede ejecutarse sobre un simulador clásico o sobre hardware
físico antes de la etapa de clasificación clásica.}
\label{fig:arquitectura}
\end{figure}

\subsection{Preprocesamiento de datos}
El dataset original supera las 70 características, incompatibles con un
procesador de 3 qubits. El preprocesamiento incluye: limpieza de
valores nulos e infinitos, conversión de etiquetas a formato binario
(benigno frente a ataque), y selección de las $k=3$ características más
discriminativas mediante \texttt{SelectKBest} (test de información
mutua). Se descartó el PCA tradicional para el modelo cuántico por
producir componentes de difícil interpretación física en espacios
reducidos. Los valores se escalan al rango $[-1,1]$ con un
\texttt{MinMaxScaler} ajustado, para evitar el fenómeno de
\textit{phase wrapping} en las compuertas de rotación del circuito.

\subsection{Arquitectura QSVM: feature map y kernel de fidelidad}
La codificación se realiza mediante un \texttt{ZZFeatureMap} con
entrelazamiento lineal. Para un vector de entrada $\mathbf{x}$, el
estado cuántico se prepara como:
\begin{equation}
  |\phi(\mathbf{x})\rangle = U_{\mathrm{FM}}(\mathbf{x})|0\rangle^{\otimes n},
\end{equation}
donde $U_{\mathrm{FM}}$ aplica rotaciones $Z$ dependientes de los datos
y compuertas $ZZ$ que entrelazan qubits adyacentes. El kernel cuántico
de fidelidad entre dos muestras se define como:
\begin{equation}
  K(\mathbf{x}_i,\mathbf{x}_j) =
  \big|\langle \phi(\mathbf{x}_i)\,|\,\phi(\mathbf{x}_j)\rangle\big|^{2}
  = \big|\langle 0|U_{\mathrm{FM}}^\dagger(\mathbf{x}_j)
    U_{\mathrm{FM}}(\mathbf{x}_i)|0\rangle\big|^{2}.
  \label{eq:kernel}
\end{equation}
Esta matriz de similitudes alimenta una SVM clásica, que aprende la
frontera de separación entre las clases. El entrelazamiento lineal (en
lugar de completo) reduce el número de compuertas CNOT cruzadas,
mitigando la decoherencia en hardware físico; en el procesador real se
emplea una profundidad de \texttt{reps}$=1$ para minimizar el tiempo de
exposición al ruido.

\subsection{Decisión de diseño: del VQC al QSVM}
El modelo inicial de esta investigación fue un VQC entrenado con
COBYLA. Sin embargo, al escalar el volumen de datos y la profundidad
del \textit{ansatz}, el entrenamiento colapsó hacia métricas cercanas
al azar, síntoma inequívoco de mesetas estériles. Los datos tabulares
de ciberseguridad carecen de la estructura geométrica que los circuitos
variacionales profundos mapean eficientemente.

La transición hacia el QSVM constituyó el punto de inflexión: al
separar el cómputo cuántico (el kernel es estático y no requiere
optimización iterativa) del entrenamiento clásico (SVM), se elimina el
riesgo de colapso por gradientes planos. Esta decisión de diseño,
documentada como lección aprendida, fue la que permitió alcanzar un
rendimiento estable y reproducible.

\subsection{Estrategia de validación en hardware}
Dada la limitación de qubits y el costo computacional del kernel, la
validación física se acota a subconjuntos reducidos de entrenamiento y
prueba. El número de circuitos de fidelidad requeridos crece según la
Tabla~\ref{tab:complejidad}.

\begin{table}[!t]
\centering
\caption{Circuitos de fidelidad requeridos por tamaño de muestra.}
\label{tab:complejidad}
\begin{tabular}{@{}ccc@{}}
\toprule
Muestras train & Muestras test & Circuitos cuánticos \\
\midrule
4  & 4  & 26 \\
6  & 4  & 45 \\
8  & 4  & 68 \\
8  & 8  & 100 \\
10 & 10 & 155 \\
\bottomrule
\end{tabular}
\end{table}

% ============================ IV. IMPLEMENTACIÓN ======================
\section{Implementación}

\subsection{Stack tecnológico}
El prototipo se desarrolló en Python, empleando Qiskit y
\texttt{qiskit-machine-learning} para el modelado cuántico
\cite{qiskit2023}; \texttt{scikit-learn}, \texttt{pandas} y
\texttt{numpy} para el preprocesamiento y el \textit{baseline} clásico;
y Streamlit para el \textit{dashboard} de visualización. La simulación
local se realizó con \texttt{qiskit-aer}.

\subsection{Generación de tráfico y detección en vivo}
Para evaluar el modelo frente a ataques realistas se implementó un
simulador de tráfico basado en Scapy, capaz de generar inundaciones TCP
SYN, UDP e ICMP, así como ataques híbridos multivector y tráfico de
fondo. Un módulo de detección en vivo captura paquetes, extrae
características por ventanas temporales y las clasifica en tiempo real.

\subsection{Integración con hardware cuántico}
La ejecución física se realiza sobre un procesador SpinQ Triangulum
(RMN, 3 qubits). La conexión se gestiona mediante la librería
\texttt{spinqit} sobre TCP/IP. A diferencia de los dispositivos
superconductores, el hardware de RMN opera mediante pulsos de
radiofrecuencia sobre un conjunto macroscópico de moléculas y devuelve
distribuciones de conteos (\textit{counts}) por cada ejecución, con
$1024$ disparos por circuito.

\subsection{Circuito de fidelidad}
La Fig.~\ref{fig:circuito} esquematiza el circuito empleado. Los tres
qubits codifican las características seleccionadas; las compuertas de
Hadamard preparan la superposición, las rotaciones $R_z$ codifican los
datos, y las compuertas CNOT introducen dependencias entre
características. La fidelidad se estima a partir de la probabilidad de
medir el estado $|000\rangle$:
\begin{equation}
  \widehat{F} \approx \frac{N_{000}}{N_{\text{shots}}},
  \label{eq:fidelidad}
\end{equation}
con $N_{\text{shots}}=1024$. El circuito prepara el estado de la
muestra A, aplica la operación inversa correspondiente a la muestra B,
y mide la probabilidad de retornar al estado fundamental.

\begin{figure}[!t]
\centering
\begin{tikzpicture}[scale=0.9]
\def\xh{0}      % etiqueta |0>
\def\xone{1.0}  % H
\def\xrz{2.0}   % Rz datos A
\def\xcnot{3.1} % CNOT
\def\xrzB{4.4}  % Rz datos B
\def\xhtwo{5.4} % H final
\def\xmeas{6.4} % medición
\def\gw{0.62}   % ancho compuerta

\foreach \y/\lbl in {2/$q_0$, 1/$q_1$, 0/$q_2$}{
  \draw[thick] (0.6,\y) -- (7.1,\y);
  \node[left, font=\scriptsize] at (\xh,\y) {$|0\rangle$};
  \node[left=2pt, font=\scriptsize] at (0.6,\y) {\lbl};
  % Medición
  \draw[fill=white] (\xmeas-0.28,\y-0.28) rectangle ++(0.56,0.56);
  \node[font=\tiny] at (\xmeas,\y) {M};
}

% Compuertas H (preparación)
\foreach \y in {2,1,0}{
  \draw[fill=blue!15] (\xone-0.28,\y-0.28) rectangle ++(0.56,0.56);
  \node[font=\scriptsize] at (\xone,\y) {H};
  % Rz (codificación A)
  \draw[fill=orange!25] (\xrz-0.28,\y-0.28) rectangle ++(0.56,0.56);
  \node[font=\tiny] at (\xrz,\y) {$R_z$};
  % Rz (codificación B, inversa)
  \draw[fill=orange!25] (\xrzB-0.28,\y-0.28) rectangle ++(0.56,0.56);
  \node[font=\tiny] at (\xrzB,\y) {$R_z^\dagger$};
  % H final
  \draw[fill=blue!15] (\xhtwo-0.28,\y-0.28) rectangle ++(0.56,0.56);
  \node[font=\scriptsize] at (\xhtwo,\y) {H};
}

% CNOT q0 -> q1
\draw[thick] (\xcnot,2) -- (\xcnot,1);
\filldraw (\xcnot,2) circle (2.2pt);
\draw (\xcnot,1) circle (0.22); \draw (\xcnot-0.22,1) -- ++(0.44,0);
\draw (\xcnot,1-0.22) -- ++(0,0.44);
% CNOT q1 -> q2
\draw[thick] (\xcnot,1) -- (\xcnot,0);
\filldraw (\xcnot,1) circle (2.2pt);
\draw (\xcnot,0) circle (0.22); \draw (\xcnot-0.22,0) -- ++(0.44,0);
\draw (\xcnot,0-0.22) -- ++(0,0.44);

\node[font=\scriptsize\itshape] at (3.85,-0.85)
  {Codificación $U(\mathbf{x}_A)$ y comparación $U^\dagger(\mathbf{x}_B)$};
\end{tikzpicture}
\caption{Esquema del circuito de fidelidad para 3 qubits. Compara dos
registros de red midiendo la probabilidad de retornar al estado
$|000\rangle$.}
\label{fig:circuito}
\end{figure}

% ============================ V. RESULTADOS ===========================
\section{Resultados Experimentales}

\subsection{Configuración experimental}
La Tabla~\ref{tab:config} resume la configuración del modelo cuántico.

\begin{table}[!t]
\centering
\caption{Configuración del modelo QSVM.}
\label{tab:config}
\begin{tabular}{@{}ll@{}}
\toprule
Parámetro & Valor \\
\midrule
Qubits & 3 \\
Feature Map & ZZFeatureMap (entrelazamiento lineal) \\
Profundidad (hardware) & \texttt{reps}$=1$ \\
Selección de características & SelectKBest ($k=3$) \\
Escalado & MinMax $[-1,1]$ \\
Kernel & FidelityQuantumKernel \\
Clasificador & SVM clásica \\
Disparos (\textit{shots}) & 1024 \\
\bottomrule
\end{tabular}
\end{table}

\subsection{Baseline clásico y QSVM en simulación}
La Tabla~\ref{tab:resultados} consolida las métricas obtenidas en los
distintos entornos evaluados.

\begin{table}[!t]
\centering
\caption{Comparación de métricas por entorno (\%).}
\label{tab:resultados}
\begin{tabular}{@{}lcccc@{}}
\toprule
Modelo / Entorno & Acc. & Prec. & Recall & F1 \\
\midrule
Random Forest (clásico) & \pendiente{completar} & \pendiente{completar}
  & \pendiente{completar} & \pendiente{completar} \\
QSVM (simulador) & 83{,}3 & \pendiente{completar}
  & \pendiente{completar} & \pendiente{completar} \\
QSVM (SpinQ, laboratorio) & 25{,}0 & 33{,}3 & 50{,}0 & 40{,}0 \\
QSVM (SpinQ, en vivo) & 75{,}0 & 100{,}0 & 66{,}7 & 80{,}0 \\
\bottomrule
\end{tabular}
\end{table}

El modelo QSVM en simulación alcanzó una exactitud del $83{,}3\%$, un
resultado significativamente superior al del VQC, que no logró
convergencia estable sobre los mismos datos. Este contraste constituye
la principal evidencia empírica que justificó la transición de
arquitectura.

\subsection{QSVM en hardware cuántico real}
La primera ejecución en el laboratorio (dataset estático) arrojó
métricas cercanas al azar ($25\%$ de exactitud), atribuibles a la
compresión extrema en 3 qubits, la decoherencia y una decodificación
heurística simplificada.

Tras refinar la selección de características, el embedding y el
criterio de decodificación, la validación sobre tráfico en vivo alcanzó
una exactitud del $75\%$, con una precisión del $100\%$ (ausencia total
de falsos positivos) y un F1-Score del $80\%$. En una validación
acotada final (4 muestras de entrenamiento y 4 de prueba) el QSVM
clasificó correctamente los 8 registros, obteniendo una matriz de
confusión con 2 verdaderos positivos y 2 verdaderos negativos; no
obstante, el kernel físico presentó una desviación media del
$12{,}94\%$ (desviación máxima del $47{,}95\%$) respecto del kernel
ideal del simulador.

% ============================ VI. DISCUSIÓN ===========================
\section{Discusión}

\subsection{Robustez del enfoque de kernel}
La transición del VQC al QSVM constituyó el hallazgo metodológico más
relevante. Los datos tabulares de ciberseguridad carecen de la
estructura geométrica que los circuitos variacionales profundos mapean
eficientemente; el enfoque de kernel, en cambio, aprovecha la
expresividad del espacio de Hilbert sin depender de bucles de
optimización cuántica, evitando las mesetas estériles. Esta lección es
transferible a otros problemas de clasificación tabular en la era NISQ.

\subsection{Impacto del ruido NISQ}
El ruido de decoherencia ($T_1$, $T_2$) y las imperfecciones de los
pulsos de radiofrecuencia introducen estados espurios en las mediciones
(p.~ej., los conteos $|01\rangle$ y $|10\rangle$ en un estado de Bell).
Curiosamente, en el experimento en vivo el hardware superó a la primera
prueba de laboratorio, fenómeno atribuible a una mayor separabilidad de
las clases en las ventanas de tráfico atacado y, posiblemente, a un
efecto de regularización por ruido. Este resultado, si bien prometedor,
debe interpretarse con cautela dado el tamaño muestral reducido.

\subsection{Limitaciones}
Las principales limitaciones son: (i) la compresión a 3 qubits, que
descarta información discriminativa; (ii) el tamaño acotado de las
muestras de prueba en hardware; y (iii) el costo cuadrático del cálculo
del kernel. Estos factores impiden, por el momento, generalizar los
resultados a escala industrial.

% ============================ VII. CONCLUSIONES =======================
\section{Conclusiones y Trabajo Futuro}
Este trabajo demuestra la viabilidad funcional de un clasificador QSVM
para la detección de anomalías en tráfico de red, tanto en simulación
como en hardware cuántico real de RMN. El QSVM superó al VQC en
estabilidad y precisión, y el hardware físico produjo clasificaciones
significativamente superiores al azar pese al ruido inherente.

Como trabajo futuro se prevé: (i) ampliar el tamaño de las muestras de
validación física; (ii) explorar embeddings con entrelazamiento
dependiente de los datos para enriquecer el kernel; (iii) aplicar
técnicas de mitigación de errores (extrapolación a ruido cero); (iv)
extender la comparación a plataformas superconductoras y de iones
atrapados vía IBM Quantum y Amazon Braket; y (v) consolidar el
\textit{dashboard} interactivo como herramienta de referencia.

% ============================ REFERENCIAS =============================
\begin{thebibliography}{9}

\bibitem{biamonte2017}
J. Biamonte, P. Wittek, N. Pancotti, P. Rebentrost, N. Wiebe y S. Lloyd,
``Quantum machine learning,'' \textit{Nature}, vol. 549, n.º 7671,
pp. 195--202, 2017.

\bibitem{preskill2018}
J. Preskill, ``Quantum Computing in the NISQ era and beyond,''
\textit{Quantum}, vol. 2, p. 79, 2018.

\bibitem{sharafaldin2018}
I. Sharafaldin, A. H. Lashkari y A. Ghorbani, ``Toward generating a new
intrusion detection dataset and intrusion traffic characterization,''
en \textit{Proc. ICISSP}, 2018, pp. 108--116.

\bibitem{qiskit2023}
Qiskit contributors, ``Qiskit: An open-source framework for quantum
computing,'' IBM Quantum, 2023.

\bibitem{ibm2024}
IBM Quantum, ``IBM Quantum Learning.'' [En línea]. Disponible:
\url{https://learning.quantum.ibm.com/}.

\end{thebibliography}

\end{document}